\chapter{Introdução}\label{cap:introducao}

A dissipação de calor em dispositivos eletrônicos compactos motiva o
desenvolvimento de tecnologias de resfriamento eficientes. O escoamento
bifásico com mudança de fase em microcanais é uma alternativa para essa
aplicação: combina elevados coeficientes de transferência de calor com pequenas
dimensões físicas \cite{cheng2023}, e o fluido refrigerante absorve grandes quantidades de
calor por unidade de área enquanto percorre o canal e passa do estado líquido para o vapor.

Para o controle e a otimização desses sistemas, é necessário identificar o padrão de
escoamento, isto é, a configuração geométrica assumida pelas fases líquida e vapor dentro
do canal. Padrões como bolha dispersa, \textit{plug}, \textit{slug}, \textit{churn} e
anular ondulado apresentam comportamentos térmicos e hidrodinâmicos distintos, relevantes
para avaliar a transferência de calor e a queda de pressão \cite{cheng2023,deoliveira2021}. A identificação dos padrões por
inspeção visual de imagens de câmera de alta velocidade é subjetiva e depende
do observador \cite{dias2025knn}, o que motiva o desenvolvimento de métodos
automáticos de classificação.

No Laboratório de Estudos Térmicos e Fluidodinâmicos (LETEF) da UNISINOS, estudos
recentes com isobutano (R600a) em microtubo horizontal de 1\,mm mostram a viabilidade
dessa classificação automática. \citeonline{deoliveira2021} apresentam um estudo
experimental abrangente com análise por transformada \textit{Wavelet} contínua (CWT),
descrevendo as respostas espectrais de pressão diferencial e do sensor óptico
associadas aos padrões.
\citeonline{dias2025knn} e \citeonline{dias2025cnn} apresentam classificadores baseados
em imagens da seção de visualização: o modelo $k$NN atinge 87\% de acurácia média de teste
nas dez dobras, enquanto a rede CNN alcança 93\% na melhor dobra e média de
88{,}1\% nas dez dobras.
Esses resultados motivam a investigação de descritores que permitem examinar quais
características das imagens se repetem entre ensaios de um mesmo padrão.

O trabalho de \citeonline{deoliveira2021} caracteriza os padrões de escoamento
aplicando a transformada \textit{Wavelet} contínua aos sinais dos sensores no domínio do
tempo, extraindo o espectro de potência \textit{Wavelet} de cada sensor individualmente.
Neste trabalho, investiga-se uma representação diferente: as sequências de imagens são
reunidas em panoramas, dos quais se extraem sinais unidimensionais de intensidade em
função da posição em pixels. Essa construção permite comparar estruturas ao longo do
panorama, mas reúne observações feitas em instantes distintos. Por isso, o eixo espacial
reconstruído não representa uma fotografia instantânea do escoamento, e seus
comprimentos de onda não são convertidos em frequências temporais ou dimensões físicas
de bolhas. A câmera permanece necessária para a aquisição.

\citeonline{schuck2019} apresentam uma formulação alternativa para o cálculo da WCO em
que a transformada é aplicada às funções de autocorrelação e correlação cruzada dos sinais
em vez de ser aplicada diretamente a eles. Essa abordagem elimina a necessidade de
suavização espectral exigida pela definição clássica \cite{torrence1998}. A questão
investigada é se essa coerência e a forma dos espectros de potência dos sinais
reconstruídos apresentam maior semelhança entre ensaios do mesmo padrão do que entre
padrões distintos. Avalia-se essa questão de maneira descritiva, sem treinar um
classificador ou pressupor que correlação elevada demonstre identidade do padrão.

% =========================================================
\section{Objetivo Geral}
% =========================================================

O objetivo geral é investigar a aplicabilidade da coerência \textit{Wavelet} alternativa
de \citeonline{schuck2019} e da comparação dos perfis de potência como descritores dos
padrões de escoamento bifásico no arranjo experimental do LETEF, a partir de sinais de
intensidade extraídos dos panoramas, avaliando a semelhança dentro de cada padrão e entre
padrões distintos.

% =========================================================
\section{Objetivos Específicos}
% =========================================================

\begin{itemize}
  \item Revisar a literatura sobre padrões de escoamento bifásico em microcanais e sobre
  métodos de análise espectral de sinais para caracterização de padrões de escoamento.
  \item Apresentar a fundamentação teórica da transformada \textit{Wavelet} contínua, do
  espectro de potência (WPS), do espectro cruzado (WCS) e da coerência \textit{Wavelet}
  (WCO), com ênfase na formulação alternativa de \citeonline{schuck2019}.
  \item Documentar e aplicar a extração da intensidade nas linhas selecionadas dos
  panoramas, o pré-processamento dos sinais e o cálculo dos espectros e da coerência.
  \item Comparar a forma dos perfis de potência pelo coeficiente de Pearson e os mapas
  de coerência por um escore ponderado, apresentando as distribuições dentro dos padrões
  e entre padrões distintos.
  \item Examinar a qualidade das entradas e a sensibilidade dos resultados ao registro
  com ressalvas de reconstrução, delimitando o alcance das evidências obtidas.
\end{itemize}

% =========================================================
\section{Organização do Trabalho}
% =========================================================

Este trabalho está organizado em cinco capítulos, além desta introdução. O
\autoref{cap:fundamentacao} apresenta a revisão bibliográfica e as bases teóricas que
fundamentam a metodologia: fenomenologia dos padrões de escoamento bifásico em microcanais e
formulação matemática das transformadas \textit{Wavelet} contínua, do espectro de potência, do
espectro cruzado e da coerência \textit{Wavelet} nas formulações clássica e alternativa. O
\autoref{cap:metodologia} descreve a metodologia, detalhando o arranjo experimental, a
reconstrução do panorama, a sequência de processamento dos sinais espaciais e a
estratégia de avaliação do método. O \autoref{cap:resultados} relata os resultados obtidos nas
comparações intra-padrão e inter-padrão, com panoramas, campos espectrais, matrizes e
tabelas de semelhança. O \autoref{cap:discussao} interpreta esses resultados, apresenta
as tabelas de sensibilidade, compara os achados com a literatura e expõe as limitações
do trabalho. Por fim, o \autoref{cap:conclusao} sintetiza as
contribuições, discute o atendimento aos objetivos e aponta os encaminhamentos para trabalhos
futuros.
